\section{Representative Cases and Discussion}\label{sec:discussion}
\subsection{Dataset Cases as Verifiable Reasoning Examples}
The appendix presents six examples from different original DWG windows, covering minimum closure, wire grouping, unknown contact states, fault diagnosis, path order, and branch interruption (Figs.~\ref{fig:appendix1}--\ref{fig:appendix3}). They illustrate annotation semantics rather than held-out model performance.

The examples distinguish three sources of electrical difficulty. First, wire-only grouping differs from continuity through closed contacts: a model must identify which edges depend on the stated operating conditions. Second, fault diagnosis requires comparing predicted observations across hypotheses, rather than merely locating a broken-looking contact. Third, a state change can be irrelevant to the queried branch, so interruption questions require tracing the actual path instead of associating every switch with every nearby terminal. Compact evidence makes these distinctions explicit through a route, a counterexample, or a short observation signature.

\subsection{Interpreting Evidence Effects}
Answer-only supervision produces the strongest direct-answer accuracy. Under common explanation prompting, E instead achieves higher confirmed evidence and joint correctness, with answer advantages concentrated in measurement planning, diagnosis sufficiency and selected state-dependent tasks. These gains coexist with losses in minimum closure, grouping and bypass questions. Evidence supervision is therefore useful when supported justifications are required, but is not uniformly superior.

The original-DWG path example illustrates a successful explicit decomposition. Neither this case nor the aggregate prompt interaction proves causal faithfulness of generated text. Longer training targets, output alignment and evaluator artifacts remain alternative explanations. The fixed-seed study cannot isolate the causal contribution of evidence content.

\subsection{Scope, Dependence, and Annotation Limitations}
The name \dataset describes the task domain; its current empirical coverage is substation secondary wiring, not all electrical schematics. The two held-out projects and 300 partially overlapping original windows limit claims about general industrial transfer. Distinct image hashes and source partitions do not eliminate repeated engineering templates or abstractly isomorphic graphs. Larger or more diverse external sources would be needed for broad cross-domain claims.

The automatic annotation pipeline can produce self-consistent but visually incorrect references if entity extraction is wrong. Source traceability facilitates correction; it does not replace independent review. T12 has no verified parallel cases, T14 has a fixed negative answer prior, and generated question templates may permit shortcuts. These limitations motivate per-task reporting, label-prior controls, and explicit separation of diagnostic tasks. Native edits represent controlled perturbations rather than realistic fault prevalence or validated maintenance procedures.

\subsection{Implications for Schematic Understanding}
The results favor a distinction between concise decision support and evidence-backed interpretation. Compact evidence supervision improves the availability of valid justifications and selected multi-condition tasks, while answer-only supervision remains preferable for overall direct-answer accuracy. The state-invariance and visual-flip failures reveal that successful explanation generation is not sufficient for selective use of electrical conditions or robust perception of altered topology. A useful next step is evidence supervision that explicitly contrasts relevant and irrelevant changes, evaluated across additional drawing sources and repeated training seeds.
